\documentclass[10pt]{article}

% --- Page Setup ---
% \usepackage[a4paper, margin=0.5cm]{geometry}
\usepackage[a4paper, left=0.5cm, right=0.5cm, top=0.3cm, bottom=0.3cm]{geometry}
\usepackage{multicol}

% \allowdisplaybreaks

\usepackage{booktabs} % Required for professional horizontal rules (\toprule, \midrule, \bottomrule)
\usepackage{array}    % Required for better column formatting

% --- Math and Code ---
\usepackage{amsmath, amssymb}
\usepackage{listings}

% --- Spacing Adjustments ---
\usepackage{enumitem}
\usepackage{titlesec}

\usepackage{nicematrix, tikz}

\usepackage{graphicx}

\usepackage{float}

\usepackage[hidelinks]{hyperref} % hidelinks removes ugly red boxes around links

% 1. Remove space between list items
\setlist{nosep, leftmargin=*} 

% 2. Shrink space around section headers
\titlespacing*{\section}{0pt}{0pt}{0pt}
\titlespacing*{\subsection}{0pt}{1ex plus .2ex minus .2ex}{0.5ex plus .1ex}

% --- NEW: Reduce header font sizes ---
\titleformat*{\section}{\normalsize\bfseries\color{blue!75!black}}
\titleformat*{\subsection}{\small\bfseries\color{magenta}}

% 3. Add a thin line between columns for readability
\setlength{\columnseprule}{0.4pt}
\setlength{\columnsep}{1.5em}

% --- Custom Code Formatting ---
\lstset{
    basicstyle=\ttfamily\scriptsize,
    breaklines=true,
    aboveskip=2pt,
    belowskip=2pt
}

\usepackage{xcolor}
% \newcommand{\sectionbox}[1]{%
%      \colorbox{blue!75!black}{\color{white}\small\bfseries~#1~}%
%    }
%    \titleformat{\section}{}{}{0em}{\sectionbox}
%    \titleformat*{\subsection}{\normalsize\bfseries}

\raggedbottom

\setlength{\parskip}{0pt}
\setlength{\baselineskip}{0pt plus 0.2pt} 


\begin{document}

\section*{Data Transfer Mechanisms}

\begin{table}[H]
    \centering
    \footnotesize 
    \renewcommand{\arraystretch}{1.2} 
    \begin{tabular}{@{} l p{2.5cm} p{1.8cm} p{2cm} p{3.2cm} p{2cm} @{}}
        \toprule
        \textbf{Mechanism} & \textbf{Controller / Nature} & \textbf{Latency} & \textbf{CPU Eff.} & \textbf{Transfer Rate / Modes} & \textbf{Debugging} \\
        \midrule
        \textbf{Polling} & 
        \textbf{CPU} \newline SW-driven active check. & 
        \textbf{Very Fast} \newline 1--3 cycles. & 
        \textbf{Very Low} \newline 100\% blocked. & 
        \textbf{Low} \newline CPU fetches per byte. & 
        \textbf{Easy} \newline Deterministic. \\
        
        \textbf{Interrupt} & 
        \textbf{CPU} \newline HW trigger, SW exec. & 
        \textbf{Slower} \newline Context switch. & 
        \textbf{High} \newline Free until signaled. & 
        \textbf{Moderate} \newline CPU still moves data. & 
        \textbf{Difficult} \newline Asynchronous. \\
        
        \textbf{DMA} & 
        \textbf{DMAC} \newline Pure HW controller. & 
        \textbf{Moderate} \newline Bus arbitration. & 
        \textbf{Very High} \newline CPU fully freed. & 
        \textbf{Maximum} \newline Burst, Steal, Transparent. & 
        \textbf{Moderate} \newline Silent RAM alters. \\
        \bottomrule
    \end{tabular}
    \caption{I/O Data Transfer Mechanisms Summary}
\end{table}

Data transfer can be initiated by the CPU or the DMAC.


\subsection*{Polling (Programmed I/O)}
CPU polls a certain I/O port continuously (using software) for data or readiness of the port to perform a data transaction.\\
CPU uses 100\% of its resource to do the polling. It cannot entertain any other requests.
% \begin{center}
%     \includegraphics[width=0.15\textwidth]{images/polling.png}
% \end{center}
\clearpage
\begin{figure}[H]
    \centering
    \includegraphics[width=0.15\textwidth]{images/polling.png}
\end{figure}
\begin{itemize}
    \item Programmer has complete control over the process, easier to test and debug.
    \item Program execution of CPU held up while waiting for I/O device to get ready, inefficient use of CPU resources.
\end{itemize}
\subsection*{Interrupt Triggered}


\begin{itemize}
    \item Uses hardware: A signal from external/internal peripheral, or a change in status of some special registers notifies the CPU that some event has occur.
    \item If the CPU decide to service the interrupt, it will suspend its current program temporarily.
    \item The interrupt will come with some interrupt vector table (IVT) index.
    \item CPU look up the IVT to check the starting address of the interrupt service routine (ISR) for the corresponding interrupt.
    \item CPU then proceed to execute the ISR linked to the interrupt that was triggered (typically a short routine).
    \item Once the ISR is completed, CPU restore the saved context, returns to the interrupted routine and continue from where it had left previously.
    \item \textbf{Interrupt latency}: delay bw CPU receiving the interrupt req, and the pt of entering the ISR.
\end{itemize}
\begin{center}
    \includegraphics[width=0.1\textwidth]{images/ivt.png}
\end{center}
\begin{itemize}
    \item CPU does not need to monitor I/O device status, more efficient use of CPU resources.
    \item But more hardware interface circuitry required between I/O device and processor. Program is slightly more complex and difficult to debug.
\end{itemize}
\subsection*{Direct Memory Access (DMA)}


DMAC is a Data Bus Controller module that performs data transfer independent of CPU.
\begin{center}
    \includegraphics[width=0.4\textwidth]{images/dmac.png}
\end{center}
\begin{itemize}
    \item CPU configures some parameters for DMAC to trf data
    \begin{itemize}
        \item Src and Dest addr of data
        \item No. of bytes to trf
        \item Trigger signal telling DMAC when to start the trf.
    \end{itemize}
    \item DMAC waits till trigger signal comes in.
    \item It requests for system bus and start data trf. (Data goes into internal buffer of DMAC before transferred to dest)
    \item After trf, DMAC releases the system bus, notifies CPU of data completion via interrupt. 
\end{itemize}
3 modes of transfer:
\begin{itemize}
    \item Burst: During control of the system bus, DMAC trf multiple units of data.
    
    
    \begin{itemize}
        \item Could be entire block of data that DMAC tasked to trf, or just a subset of that block.
        \item CPU can continue with its operation as long it does not need the particular system bus that DMAC is using.
        \item So Burst has faster trf rate but CPU may be suspended/inactive if CPU has to wait for DMAC to complete its trf.
    \end{itemize}
    \item Cycle stealing: DMAC trfs one unit of data aft taking control, then releases control of system bus to CPU.
    
    
    \begin{itemize}
        \item DMAC can execute the data trf between CPU instructions or between pipeline stages.
        \item CPU may be suspended if it need to access the system bus but suspend time is shorter as only one unit is transferred at one time.
        \item Transfer rate is slower than in Burst Mode but will CPU will only be inactive for very short period of time.
        \item Favoured if CPU needs to be responsive.
        \begin{center}
            \includegraphics[width=0.4\textwidth]{images/burst-cycle-stealing.png}
        \end{center}
    \end{itemize}
    \item Transparent: DMAC trfs data only when CPU not using the system bus.
    
    
    \begin{itemize}
        \item CPU performance not affected, but potentially slowest trf rate.
        \item More complex hardware needed to detect when CPU not using the system bus.
    \end{itemize}
    \item So from burst to cycle-stealing to transparent, data trf gets slower but CPU performance/response improves.
    \item Transferring via CPU can be faster than DMA as CPU speed is typically faster, just that it means wasting CPU resources to just transfer data.
\end{itemize}

\end{document}
